\documentclass[11pt,UTF8,fontset=none]{ctexart}

\usepackage[a4paper,margin=2.3cm]{geometry}
\usepackage{fontspec}
\usepackage{xcolor}
\usepackage{listings}
\usepackage{booktabs}
\usepackage{tabularx}
\usepackage{enumitem}
\usepackage{tikz}
\usetikzlibrary{arrows.meta,positioning,fit,backgrounds,shapes.geometric,calc}
\usepackage[colorlinks=true,linkcolor=black,urlcolor=blue]{hyperref}

% --- Fonts: CJK via Noto CJK SC (verified present); Latin default Latin Modern; mono DejaVu ---
\setCJKmainfont{Noto Serif CJK SC}
\setCJKsansfont{Noto Sans CJK SC}
\setCJKmonofont{Noto Sans CJK SC}
\setmonofont{DejaVu Sans Mono}[Scale=0.88]

% --- Code listing style (ASCII source only; Chinese annotations live in prose) ---
\definecolor{framegray}{gray}{0.45}
\definecolor{bgcode}{gray}{0.97}
\lstdefinestyle{ysx}{
  basicstyle=\ttfamily\footnotesize,
  breaklines=true,
  breakatwhitespace=false,
  columns=fullflexible,
  keepspaces=true,
  showstringspaces=false,
  frame=single,
  rulecolor=\color{framegray},
  backgroundcolor=\color{bgcode},
  framesep=4pt,
  xleftmargin=4pt,
  xrightmargin=2pt,
  aboveskip=7pt,
  belowskip=7pt,
  postbreak=\mbox{\textcolor{framegray}{$\hookrightarrow$}\space},
}
\lstset{style=ysx}
\newcommand{\fcite}[1]{{\footnotesize\itshape\textcolor{framegray}{#1}}}

% Let long monospace file paths reflow instead of running into the right margin.
\makeatletter
\g@addto@macro\UrlBreaks{\do\:\do\-}
\makeatother
\DeclareUrlCommand\fp{\urlstyle{tt}}
\setlength{\emergencystretch}{3.5em}
\tolerance=1200
\hbadness=4000

\title{\bfseries 从 riscv-opcodes 到端到端指令支持\\[2pt]
\large YSX auto-td 自动 TableGen 生成技术报告}
\author{昇腾 950 CCU 编译器 · 内部技术资料}
\date{2026-06-15 · 分支 \texttt{ysx-tiny-fv-ccu}（基线 \texttt{f5fd71c})}

\begin{document}
\maketitle

\noindent\small\textit{本报告不含性能测试数据。所有代码/YAML/TableGen 片段均标注其在工程中的文件与行号区间，行号以本分支当前树为准；代码清单为 ASCII 源逐字呈现，中文注解一律置于清单外的正文，凡为版面计裁剪者均标注“节选”。}\normalsize

\section*{摘要}
\addcontentsline{toc}{section}{摘要}
我们为昇腾 950 CCU 实现的基础编译器，在 ISA 裁剪后以标量 \texttt{rv64ima} 为基线。为支撑数据加载与 reduce 等向量操作，需要在此基线上补回一小批 RVV 指令。本工作不沿用“为每条指令手写嵌套 TableGen 指令类”的传统路径，而是构建了 \textbf{auto-td} 生成系统：每条新指令由\textbf{一个结构化 YAML 清单}声明其身份，其\textbf{编码 100\% 来自 opcode 数据库}（标准指令取自外部权威 \texttt{riscv-opcodes}，自定义指令取自本团队维护的 \texttt{ysx-opcodes}）、\textbf{操作数与副作用 100\% 来自共享 taxonomy}，由 Python 生成器产出真正的指令 TableGen 记录、审计 manifest，以及（仅对显式标记 \texttt{builtin.codegen:\ true} 的证明 API）Clang builtin TD、LLVM intrinsic TD 与 CGBuiltin 派发片段。当前切片覆盖 50 条指令（15 tiny-F + 35 tiny-V，其中 49 条来自 \texttt{riscv-opcodes}、1 条 \texttt{yushuxin.vfexp} 来自 \texttt{ysx-opcodes}）。

需要诚实界定端到端的广度：50 条中\textbf{只有 8 条}（\texttt{vadd.vv}/\texttt{vsub.vv}/\texttt{vmul.vv}/\texttt{vredsum.vs}/\texttt{vfredusum.vs}/\texttt{vrgather.vv}/\texttt{vslideup.vx}/\texttt{yushuxin.vfexp}）带生成的 builtin+intrinsic+CGBuiltin 并具备从 C 到目标码的路径，其余 42 条仅到“指令记录 + MC 编码”层；且这 8 条的 intrinsic$\to$机器\textbf{下降仍是手写 C++}，用户头文件 \texttt{ysx\_vector.h} 也是手写的（生成器不产 \texttt{.h}）。

本报告的核心是一组\textbf{新旧对比}：传统上端到端引入一条向量指令需要在多个文件、约六层 TableGen 类/多类之间铺陈编码位、操作数、调度、pseudo、pattern、intrinsic、builtin；而在 auto-td 下，新增一条指令在常见情形下\textbf{只需新增一个 YAML 文件}（必要时加一个 taxonomy 类目）。我们同时对一个常被忽略的命题给出\textbf{诚实分层}判断：上游 TableGen 的“人为知识”中，\textbf{编码/操作数位置/汇编–反汇编层}确实可由 opcode 数据库机械派生，但\textbf{语义/ISel/调度/intrinsic 类型/ABI}是真正的人为设计，无法靠解析 \texttt{riscv-opcodes}（或散文式 \texttt{riscv-isa-manual}）立刻得到。auto-td 的创新不是“消灭人为知识”，而是\textbf{把机械层与设计层分离}。最后，我们补强了守护，使“新增 tiny-F/V 指令记录必须经 auto-td”这一承诺，在\textbf{其覆盖的形态内}，从约定变为构建/CI 门禁。

\section{引言}
在 LLVM 后端引入一条新指令的工程成本，主要不在“这条指令本身”，而在它必须被同时表达成多种彼此独立的事实：机器码编码、汇编/反汇编语法、操作数与寄存器类、副作用与调度、选择 DAG 模式、intrinsic 原型、前端 builtin 与头文件。传统 RISC-V 后端用一套深层 TableGen 类/多类库来复用这些表达，但代价是：\textbf{加一条指令前，先要读懂并挑对这套类层级}；遇到库里没有的新形态，还要先扩库。

本工作提出的问题是：这些表达里，有多少其实是\textbf{可由权威数据源机械派生}的，有多少是\textbf{真正的人为设计}？贡献为：(1) 一个把两者显式分离的 auto-td 生成系统；(2) 一条自定义指令 \texttt{yushuxin.vfexp} 的端到端存在性证明；(3) 一道把“无手写指令 TD”承诺机械化的源端守护；(4) 对上述命题的诚实分层判断。全文以“旧路 vs 新路”为主线（第 3–6 节）。图~\ref{fig:flow}、\ref{fig:oldnew}、\ref{fig:guard} 分别给出数据流、新旧对比与守护关系。

\section{背景：上游 RVV 指令的完整 bring-up 路径}
以一条最普通的整数向量加法 \texttt{vadd.vv} 为例，上游 LLVM/Clang 把它表达成以下几类事实（以下 TableGen 片段均为节选）。

\textbf{(a) 编码与汇编（格式类，手工誊写编码位）。} 编码位被手写进格式类的 \texttt{let Inst\{\}}。注意末行 \texttt{RVVConstraint} 是设计层信息、非编码，opcode DB 里没有：
\begin{lstlisting}
class RVInstVV<bits<6> funct6, RISCVVFormat opv, dag outs, dag ins,
               string opcodestr, string argstr>
    : RVInst<outs, ins, opcodestr, argstr, [], InstFormatR> {
  bits<5> vs2; bits<5> vs1; bits<5> vd; bit vm;
  let Inst{31-26} = funct6;     let Inst{25}    = vm;
  let Inst{24-20} = vs2;        let Inst{19-15} = vs1;
  let Inst{14-12} = opv.Value;  let Inst{11-7}  = vd;
  let Inst{6-0}   = OPC_OP_V.Value;
  let Uses = [VL, VTYPE];
  let RVVConstraint = VMConstraint;
}
\end{lstlisting}
\fcite{llvm/lib/Target/RISCV/RISCVInstrFormatsV.td:107--125（class RVInstVV，节选）}

\textbf{(b) 指令记录（嵌套多类）。} 实际指令经多类库实例化，多类之间层层 mix-in：
\begin{lstlisting}
multiclass VALU_IV_V_X_I<string opcodestr, bits<6> funct6>
    : VALU_IV_V<opcodestr, funct6>, VALU_IV_X<opcodestr, funct6>,
      VALU_IV_I<opcodestr, funct6>;
...
defm VADD_V : VALU_IV_V_X_I<"vadd", 0b000000>;
\end{lstlisting}
\fcite{llvm/lib/Target/RISCV/RISCVInstrInfoV.td:612 与 :1137}

\textbf{(c) Pseudo 与选择模式。} 另有一套 VPseudo/VPat 多类，规模庞大（\texttt{RISCVInstrInfoVPseudos.td} 共 7332 行）；一条 \texttt{defm} 会展开出\textbf{完整的 LMUL/掩码/policy 矩阵}：
\begin{lstlisting}
multiclass VPseudoBinaryV_VV<LMULInfo m, ...> { ... }
multiclass VPatBinaryV_VV<string intrinsic, string instruction, ...> { ... }
...
defm PseudoVADD : VPseudoVALU_VV_VX_VI<Commutable=1>;
\end{lstlisting}
\fcite{llvm/lib/Target/RISCV/RISCVInstrInfoVPseudos.td:2160 / :4810 / :6111}

\textbf{(d) LLVM intrinsic} 由 \texttt{defm vadd : RISCVBinaryAAX;}（\texttt{IntrinsicsRISCV.td:1199,1371}）声明；\textbf{(e) 前端 builtin} 由 \texttt{defm vadd : RVVIntBinBuiltinSet;}（\texttt{riscv\_vector.td:759}）声明。

值得强调：上游\textbf{前端这层本身已经是生成的}。\texttt{clang/utils/TableGen/RISCVVEmitter.cpp} 中的 \texttt{RVVEmitter} 后端从 \texttt{riscv\_vector.td} 出发，分别产出 builtin 表、CodeGen 派发、sema 与 \texttt{riscv\_vector.h} 头文件（\texttt{RVVEmitter::createHeader}:407、\texttt{createBuiltins}:505、\texttt{createCodeGen}:573）。也就是说，上游真正的“人为知识”集中在\textbf{后端}：格式类的 \texttt{let Inst\{\}}、以及 VALU/VPseudo/VPat 这套深层多类库。auto-td 的发力点正是这一层（前端侧两者都已生成，auto-td 在 Clang 侧并不更优，见 §5）。

\section{旧方法基线（对比的一侧）}
把上节归纳成“端到端加一条 \texttt{vadd.vv} 需要做什么”：(1) 选/写格式类并把编码位手工誊进 \texttt{let Inst\{\}}（\texttt{RISCVInstrFormatsV.td}）；(2) 选对 ALU 多类并 \texttt{defm} 出指令记录（\texttt{RISCVInstrInfoV.td}）；(3) \texttt{defm} 出 VPseudo 与 VPat（\texttt{RISCVInstrInfoVPseudos.td}）；(4) \texttt{defm} 出 LLVM intrinsic（\texttt{IntrinsicsRISCV.td}）；(5) \texttt{defm} 出 Clang builtin（\texttt{riscv\_vector.td}），其后前端四件套由 \texttt{RVVEmitter} 生成；(6) MC 编/解码由格式类的 \texttt{let Inst\{\}} 自动派生（上游已自动化）。

常见情形约触及 4–5 个 \texttt{defm}，\textbf{但其真实成本是这些 \texttt{defm} 背后约六层类/多类的间接}：作者必须先理解 \texttt{VALU\_IV\_V\_X\_I} $\to$ \texttt{VALU\_IV\_V/X/I} $\to$ \texttt{VALUVV} $\to$ \texttt{RVInstVV} 与 \texttt{VPseudo*/VPat*} 的层级，才能挑对复用点；一旦是库里没有的\textbf{新形态}（例如一条自定义逐元素函数指令），就得\textbf{新增/扩展格式类与多类机制}（格式类 \texttt{.td}、ALU 多类 \texttt{.td}、\texttt{VPseudos.td}、\texttt{VPat}、\texttt{IntrinsicsRISCV.td}、\texttt{riscv\_vector.td}，乃至 per-CPU 调度模型 \texttt{RISCVSched*.td}），跨多个文件，文件与心智成本显著上升。换言之：编码位被人手誊写、形态被人手挑类——而这两件事，恰恰是最“机械”、最可由数据派生的部分。

\section{YSX auto-td 系统设计（对比的另一侧）}

\subsection{权威数据分层}
auto-td 把“加一条指令”拆成三种各有归属的事实。\textbf{opcode 数据库拥有编码位}：变量字段位置取自 \texttt{arg\_lut.csv}，固定位取自 opcode 行内的 \texttt{msb..lsb=val} 语法。解析逻辑只做“把无等号的 token 当变量字段、有等号的当固定位”：
\begin{lstlisting}
fields = tuple(part for part in parts[1:] if "=" not in part)
fixed_bits = tuple(part for part in parts[1:] if "=" in part)
field_ranges = tuple(OpcodeFieldRange(field, *arg_lut[field])
                     for field in fields if field in arg_lut)
records[key] = OpcodeRecord(key, mnemonic, fields, fixed_bits, path, field_ranges)
\end{lstlisting}
\fcite{llvm/lib/Target/YuShuXin/auto-td/tools/ysx\_auto\_td/opcodes.py:23--32}

\textbf{taxonomy 拥有共享的操作数/副作用形态}（按类目，而非按指令）：它引用 opcode 的字段\emph{名}（\texttt{vd}/\texttt{vs1}/\texttt{vs2}/\texttt{vm}）作为连接键，赋予角色/寄存器类/副作用——这些是 opcode 里\emph{没有}的设计信息。下列为节选（省略 \texttt{domain}、\texttt{default\_surfaces} 与各 \texttt{false} 副作用行）：
\begin{lstlisting}
custom_float_unary:
  operation: fexp
  operands:
    outs: [{role: dest, field: vd, reg_class: VR}]
    ins:
      - {role: value, field: vs2, reg_class: VR}
      - {role: mask_policy, field: vm, operand: VMaskOp}
  effects: {implicit_uses: [VL, VTYPE]}
  sched: {semantic_class: vector_float_alu}
\end{lstlisting}
\fcite{llvm/lib/Target/YuShuXin/auto-td/taxonomy/vector-alu.yaml:66--80（custom\_float\_unary，节选）。其中 \texttt{sched} 当前仅信息性，生成器恒置 \texttt{hasNoSchedulingInfo=1}。}

\textbf{一个 YAML 拥有指令身份}（助记符、opcode 来源、\texttt{spec\_ref}、特性要求，及可选的 pseudo/pattern/builtin 事实），且\emph{不含任何编码位}（见 §4.3 禁词守护）。注意 \texttt{spec\_ref}（如 \texttt{tinyv.vector-alu.custom\_float\_unary}）只是 \textbf{taxonomy 类目的索引键}（\texttt{loader.py:184--185} 取 \texttt{rsplit('.')[-1]}），它\emph{不指向} \texttt{riscv-isa-manual}；\texttt{isa-manual} 当前未被任何工具或字段机器引用（见 §9）。

\subsection{生成扁平记录（无逐指令嵌套类）}
生成器把上述事实汇成一条\textbf{单一共享基类 \texttt{RVInst} 的扁平 \texttt{def}}——没有任何逐指令的派生类层级：
\begin{lstlisting}
def _emit_instruction(instruction) -> list[str]:
    record_name = _record_name(instruction)        # "YSX_AUTO_" + 规范化助记符
    outs, ins, asm_operands = _operand_dags(instruction)   # 来自 taxonomy
    assignments = _assignments(instruction)         # 来自 opcode 编码位
    lines = [ f'def {record_name} : RVInst<{outs}, {ins}, "{mnem}", ...> {{',
              f'  let Predicates = [{_predicate(instruction)}];', ... ]
    # 恒发 hasSideEffects/mayLoad/mayStore/hasNoSchedulingInfo，再逐位写 let Inst{...}
def _record_name(instruction):
    return "YSX_AUTO_" + source_key_from_mnemonic(instruction.mnemonic).upper()
\end{lstlisting}
\fcite{llvm/lib/Target/YuShuXin/auto-td/tools/ysx\_auto\_td/emit\_td.py:334--366（节选）}

数据流与三处 CMake fan-out 见图~\ref{fig:flow}。

\begin{figure}[ht]
\centering
\begin{tikzpicture}[
  font=\footnotesize,
  node distance=6mm and 9mm,
  src/.style={draw,rounded corners=2pt,fill=blue!6,minimum height=7mm,minimum width=20mm,align=center},
  stage/.style={draw,fill=gray!8,minimum height=7mm,minimum width=16mm,align=center},
  tree/.style={draw,rounded corners=2pt,fill=green!7,minimum height=10mm,minimum width=33mm,align=center},
  ar/.style={-{Latex[length=2mm]}}]
  \node[src] (yaml) {指令 YAML};
  \node[src,below=3mm of yaml] (tax) {taxonomy};
  \node[src,below=3mm of tax] (op) {opcode DB};
  \node[stage,right=of tax] (loader) {loader};
  \node[stage,right=of loader] (val) {validate};
  \node[stage,right=of val] (emit) {emit\_td};
  \node[stage,right=of emit] (rep) {report};
  \draw[ar] (yaml.east) -- (loader.west);
  \draw[ar] (tax.east) -- (loader.west);
  \draw[ar] (op.east) -- (loader.west);
  \draw[ar] (loader) -- (val);
  \draw[ar] (val) -- (emit);
  \draw[ar] (emit) -- (rep);
  \node[tree,below=10mm of loader] (t1) {LLVM target 树\\\scriptsize YSXInstrInfo.td 含 *.inc};
  \node[tree,below=10mm of emit] (t2) {LLVM IR 树\\\scriptsize IntrinsicsYSX.td 含 *Intrinsics.td};
  \node[tree,right=8mm of t2] (t3) {Clang 树\\\scriptsize BuiltinsYSX.td / RISCV.cpp};
  \draw[ar] (emit.south) .. controls +(0,-6mm) and +(0,6mm) .. (t1.north);
  \draw[ar] (emit.south) -- (t2.north);
  \draw[ar] (emit.south) .. controls +(0,-6mm) and +(0,6mm) .. (t3.north);
\end{tikzpicture}
\caption{auto-td 数据流与三处 CMake fan-out：指令 YAML + taxonomy + opcode DB 经 loader/validate/emit\_td/report，扇出到 LLVM target、LLVM IR、Clang 三棵构建树。}
\label{fig:flow}
\end{figure}

\subsection{守护：禁止把“已派生”的事实手写回 YAML}
校验器禁止在指令 YAML 里夹带原始 TableGen 或拷贝来的编码事实，从源头保证“机械层只来自 opcode”：
\begin{lstlisting}
FORBIDDEN = ("raw_td", "def : Pat", "RVInst", "VUnitStrideLoad",
             "VRED_", "VPseudo", "let Inst{", "bits<")
FORBIDDEN_FIELDS = {"raw_cpp", "encoding", "fixed_bits"}
\end{lstlisting}
\fcite{llvm/lib/Target/YuShuXin/auto-td/tools/ysx\_auto\_td/validate.py:4--14}

\section{新旧对比分析（核心章）}
图~\ref{fig:oldnew} 给出旧路深层多类间接与新路扁平生成的对照；表~\ref{tab:cmp} 逐维展开。

\begin{figure}[ht]
\centering
\begin{tikzpicture}[font=\scriptsize,
  b/.style={draw,rounded corners=1.5pt,minimum width=42mm,minimum height=5.6mm,align=center,inner sep=2pt},
  old/.style={b,fill=red!6}, new/.style={b,fill=green!8},
  hd/.style={align=center,font=\footnotesize\bfseries},
  ar/.style={-{Latex[length=1.7mm]}}]
  % ---- left column (old), centered at x=0 ----
  \node[hd] (oh) at (0,0) {旧路：约 6 层类/多类间接};
  \node[old,below=2.4mm of oh] (o1) {RVInstVV（格式类，手写 let Inst）};
  \node[old,below=3mm of o1] (o2) {VALUVV};
  \node[old,below=3mm of o2] (o3) {VALU\_IV\_V / \_X / \_I};
  \node[old,below=3mm of o3] (o4) {VALU\_IV\_V\_X\_I};
  \node[old,below=3mm of o4] (o5) {defm VADD\_V};
  \node[old,below=3mm of o5,fill=red!12] (o6) {+ VPseudo*/VPat* + Intrinsics\\+ riscv\_vector};
  \draw[ar] (o1)--(o2); \draw[ar](o2)--(o3); \draw[ar](o3)--(o4); \draw[ar](o4)--(o5); \draw[ar](o5)--(o6);
  % ---- right column (new), centered at x=6.6cm (no overlap: 42mm boxes, 24mm gap) ----
  \node[hd] (nh) at (6.6,0) {新路：一个 YAML $\to$ 生成};
  \node[new,below=2.4mm of nh] (n1) {指令 YAML + taxonomy\\+ opcode DB（三类输入）};
  \node[new,below=6mm of n1] (n2) {ysx\_auto\_td\_gen.py};
  \node[new,below=6mm of n2] (n3) {扁平 def YSX\_AUTO\_*\\（一条记录，无嵌套类）};
  \draw[ar] (n1)--(n2); \draw[ar] (n2)--(n3);
\end{tikzpicture}
\caption{旧路 vs 新路。旧路：作者须沿约六层类/多类挑对复用点；新路：复杂度集中在生成器，单指令呈一条扁平 \texttt{def}。}
\label{fig:oldnew}
\end{figure}

\begin{table}[ht]
\centering\footnotesize
\renewcommand{\arraystretch}{1.25}
\begin{tabularx}{\linewidth}{@{}>{\raggedright\arraybackslash}p{0.16\linewidth} X X@{}}
\toprule
\textbf{维度} & \textbf{旧路（上游 RVV bring-up）} & \textbf{新路（YSX auto-td）}\\
\midrule
触及文件 & 常见约 4–5 个 \texttt{defm}（背后约六层类/多类）；新形态需跨多文件新增/扩展格式类与多类 & \textbf{一个 YAML}；新形态再加一个 taxonomy 类目\\
TableGen 类机制 & 约 6 层类/多类间接（格式类$\to$ALU 多类$\to$VPseudo$\to$VPat） & \textbf{零}逐指令类；统一 \texttt{RVInst} 扁平 \texttt{def}\\
编码位来源 & 手写进格式类 \texttt{let Inst\{\}} & 100\% 派生自 opcode DB，YAML 禁含编码\\
操作数/副作用 & 手写进操作数类与 \texttt{let mayLoad/...} & 100\% 来自 taxonomy 类目（按字段名连接）\\
pseudo/pattern 矩阵$^\dagger$ & 单条 \texttt{defm} 展开\textbf{完整 LMUL/掩码/policy 矩阵} & \textbf{不生成}：\texttt{*Pseudos/*Patterns.inc} 为注释 manifest\\
intrinsic/builtin 胶水 & \texttt{defm} + \texttt{RVVEmitter} 生成前端四件套（含头文件） & \texttt{codegen:true} 时生成 builtin TD/intrinsic TD/CGBuiltin；\textbf{不含用户头文件}\\
MC 层 & 自 \texttt{let Inst\{\}} 自动派生 & 自生成记录 \texttt{let Inst\{\}} 自动派生（同机制）\\
ISel/lowering & 标准 intrinsic 走声明式 VPat & \textbf{手写 C++}（全部 8 条 codegen 指令，见 §6）\\
评审面 & 跨多文件读 \texttt{defm}，核对挑类 & 读一个约 10 行 YAML（+ 可能一个类目）\\
复杂度所在 & 深层手维护的多类库 & 集中在一个 Python 生成器 + 声明式 taxonomy\\
不变量强制 & N/A（手写即正道） & 源端守护 + 测试注册门禁（§7，限其覆盖形态）\\
\bottomrule
\end{tabularx}
\caption{新旧逐维对比。}
\label{tab:cmp}
\end{table}

\noindent$^\dagger$ 这是对比中最需诚实的一格：上游一条 \texttt{defm PseudoVADD} 会展开出完整的伪指令与模式矩阵，是真正的功能产物；auto-td 的扁平 \texttt{def} 每条只生成一条指令记录，\texttt{YSXGenAutoTinyVPseudos.inc} 当前为空（\texttt{no generated records yet}），相应矩阵尚未下降——因此本格不是“新路更省”，而是“新路尚未覆盖”。

\medskip
\noindent\textbf{命题的诚实分层判断。}“上游 TableGen 的人为知识本可由 opcodes/isa-manual 解析获取”——
\begin{itemize}[leftmargin=1.4em,itemsep=2pt]
\item 对\textbf{机械层成立}：\texttt{RVInstVV} 里那串 \texttt{let Inst\{31-26\}=funct6 ... 6-0=OPC\_OP\_V}（\texttt{RISCVInstrFormatsV.td:107--125}）所誊写的，与 opcode 行 \texttt{vadd.vv 31..26=0x00 vm vs2 vs1 14..12=0x0 vd 6..0=0x57} 是\emph{同一份信息}。编码、操作数字段位置、汇编/反汇编，确实可由 opcode DB 立刻派生——auto-td 正是这么做的。
\item 对\textbf{语义/设计层不成立}：操作数\emph{角色}与寄存器类、\texttt{mayLoad/mayStore}、隐式 \texttt{VL/VTYPE}、ISel 语义、调度、intrinsic 类型、ABI——\texttt{riscv-opcodes} 里\emph{没有}（连 \texttt{RVInstVV} 都还要手写 \texttt{RVVConstraint}），\texttt{riscv-isa-manual} 只有自然语言散文，无法靠解析立刻得到。
\item 因此 auto-td 的创新是\textbf{分离}：机械层完全派生；设计层收敛为 taxonomy 这份\textbf{集中、去重的残余人为知识}（opcodes 给编码、taxonomy 给共享形态、YAML 给身份）+ 极少量算法性手写 C++。这份残余知识在\textbf{指令共享形态时是次线性的}（49 条 RVV 指令复用少数类目）；但对\textbf{形态各异的厂商自定义指令}，每条约需一个新 taxonomy 类目 + 手写下降，\textbf{近似按条计}（如 \texttt{custom\_float\_unary} 仅服务 \texttt{vfexp} 一条）。
\end{itemize}

\section{端到端案例：yushuxin.vfexp}
\texttt{yushuxin.vfexp} 是本切片中唯一的厂商自定义指令，用作“只往数据源加必要信息即可端到端”的\textbf{存在性证明}（非代表全部 50 条；其余 42 条只到记录+MC 层）。

\textbf{第一步，opcode 一行}（与 \texttt{riscv-opcodes} 同语法；注意此行由本团队写入 \texttt{ysx-opcodes}，是“单源录入”而非外部权威）：
\begin{lstlisting}
yushuxin.vfexp 31..26=0x2a vm vs2 19..15=0 14..12=0x1 vd 6..0=0x0b
\end{lstlisting}
\fcite{third\_party/ysx-opcodes/extensions/rv\_xtinyv:1}

\textbf{第二步，一个 YAML 清单}（无任何编码位）。其 \texttt{header:} 字段是审计元数据、\emph{不生成指令}（生成器不产 \texttt{.h}，用户头文件手写，见第六步）；\texttt{masked:\ true} 与多 LMUL 当前\emph{尚未实现}（ISel 仅下降 \texttt{v4f32} 非掩码路径）：
\begin{lstlisting}
mnemonic: yushuxin.vfexp
opcode_source: {repo: ysx-opcodes, extension: rv_xtinyv, key: yushuxin_vfexp}
spec_ref: tinyv.vector-alu.custom_float_unary
features: {required: [xtinyv]}
pseudos:
  matrix: {element_types: [f32], lmuls: standard, masked: true, policy: llvm_default}
patterns:
  - {kind: intrinsic_to_pseudo, intrinsic: ysx.vfexp, operation: fexp}
builtin:
  header: ysx_vector.h
  names: [ysx_vfexp_v_f32m1]
  overloaded: false
  codegen: true
\end{lstlisting}
\fcite{llvm/lib/Target/YuShuXin/auto-td/instructions/tiny-v/yushuxin\_vfexp.yaml:1--13}

\textbf{第三步，运行生成器}，得到真正的指令记录（由 \texttt{emit\_td.py:\_emit\_instruction} 产出，记录名 \texttt{YSX\_AUTO\_YUSHUXIN\_VFEXP}；落在构建树 \texttt{YSXGenAutoTinyVInstrInfo.inc}，committed 树中不存在任何手写版本；构建树行号随树而异，故以记录名引用）。以下为\textbf{逐字}生成内容：
\begin{lstlisting}
// opcode-source: ysx-opcodes/rv_xtinyv/yushuxin_vfexp
def YSX_AUTO_YUSHUXIN_VFEXP : RVInst<(outs VR:$vd), (ins VR:$vs2, YSXAutoVMaskOp:$vm), "yushuxin.vfexp", "$vd, $vs2$vm", [], InstFormatR> {
  let Predicates = [HasStdExtXTinyV];
  let hasSideEffects = 0;
  let mayLoad = 0;
  let mayStore = 0;
  let hasNoSchedulingInfo = 1;
  let Uses = [VL, VTYPE];
  bit vm;
  bits<5> vs2;
  bits<5> vd;
  let Inst{31-26} = 0b101010;
  let Inst{25} = vm;
  let Inst{24-20} = vs2;
  let Inst{19-15} = 0b00000;
  let Inst{14-12} = 0b001;
  let Inst{11-7} = vd;
  let Inst{6-0} = 0b0001011;
}
\end{lstlisting}
\fcite{由 emit\_td.py 逐字生成；逐位对照：\texttt{Inst\{31-26\}=0b101010} 即 opcode 行 \texttt{31..26=0x2a}，\texttt{14-12=0b001} 即 \texttt{14..12=0x1}，\texttt{6-0=0b0001011} 即 \texttt{6..0=0x0b}——编码位机械派生，非手写。}

对 \texttt{builtin.codegen:\ true}，生成器还产出 Clang builtin（\texttt{vfexp\_v\_f32m1}）、LLVM intrinsic（\texttt{int\_ysx\_vfexp}）与 CGBuiltin 派发，分别被手写薄封装 \texttt{BuiltinsYSX.td} / \texttt{IntrinsicsYSX.td} / \texttt{RISCV.cpp} 引入。\textbf{它替代了第 3 节的步骤 1–3}（选/写格式类誊编码、为新形态新增多类与 \texttt{defm}）全部被“opcode 一行 + YAML 一份 + 跑生成器”取代；步骤 4–5 的 intrinsic/builtin \emph{TD} 亦由 \texttt{codegen:true} 自动产出。

\textbf{端到端测试链。} \texttt{ysx\_vfexp\_v\_f32m1}（手写头文件 \texttt{ysx\_vector.h}）$\to$ \texttt{\_\_builtin\_ysx\_vfexp\_v\_f32m1} $\to$ \texttt{llvm.ysx.vfexp} $\to$ 生成记录。两个测试分别守住不同层：\texttt{yushuxin-vfexp.c} 守 \textbf{C$\to$IR}（断言 \texttt{call <4 x float> @llvm.ysx.vfexp}）；\texttt{tinyv-builtins-asm.c} 才是\textbf{目标码}证明（\texttt{-emit-obj} + \texttt{llvm-objdump}，\texttt{// OBJ:\ yushuxin.vfexp} 于该文件第 100 行）。

\textbf{诚实的 gap（系统性，非 vfexp 个例）。} intrinsic$\to$机器\textbf{下降对全部 8 条 \texttt{codegen} 指令均为手写 C++}，各自硬限于 \texttt{v4i32}/\texttt{v4f32} 固定形态，且\textbf{当前无任何生成的 DAG pattern}：
\begin{lstlisting}
case Intrinsic::ysx_vadd: { ... }
case Intrinsic::ysx_vsub: case Intrinsic::ysx_vmul:
case Intrinsic::ysx_vredsum: case Intrinsic::ysx_vfredsum:
case Intrinsic::ysx_vrgather: {
  if (...) { if (VT != MVT::v4f32) ... } else if (VT != MVT::v4i32) { ... }
}
case Intrinsic::ysx_vslideup: { if (... || VT != MVT::v4i32) break; ... }
case Intrinsic::ysx_vfexp:   { if (... || VT != MVT::v4f32) break; ... }   // :1195
\end{lstlisting}
\fcite{llvm/lib/Target/YuShuXin/YSXISelDAGToDAG.cpp:1133--1206（8 条 codegen 指令的手写下降，节选）}

\texttt{*Pseudos/*Patterns.inc} 是注释式审计 manifest、不是 DAG pattern；真正的下降是上面这段手写 C++。ISel 硬限固定向量类型，故 \texttt{masked:true} 与多 LMUL 当前\textbf{是愿景，尚未下降}。用户头文件 \texttt{clang/lib/Headers/ysx\_vector.h}（8 个手写 inline 封装）\textbf{整份手写}。指令记录、MC 编码、intrinsic TD、Clang builtin TD 是真实且经测试的；剩余的 \texttt{vsetvli} 插入、intrinsic 选择与公共头文件是有界的手写部分。

\section{守护与威胁分析}
审计发现一个 HIGH 级缺口：承诺“新增指令必须经 auto-td”此前\textbf{仅靠约定}——\texttt{validate.py} 只读 YAML、从不读 \texttt{.td}，且整套单测只在手动 \texttt{unittest} 下运行，不门禁构建或 CI。据此补强（关系见图~\ref{fig:guard}）。

\textbf{(1) 源端守护}：新增 \texttt{test\_no\_handwritten\_instruction\_td.py}，扫描 committed \textbf{顶层}后端 \texttt{.td}，对 \texttt{def\ ...\ :\ RVInst*} 形态的指令记录，若 body 带 tiny-F/V 谓词或 tiny-FV 助记符且落在生成 include 之外即失败，并断言“生成的 \texttt{YSX\_AUTO\_*} 记录集合 == YAML 声明的助记符集合”：
\begin{lstlisting}
def _instruction_records(td_path): ...     # 解析单行/块状两种 def (66-99)
for name, lineno, body in _instruction_records(td_path):      # 判定循环 105-112
    has_predicate = any(pred in body for pred in TINY_FV_PREDICATES)
    mnemonic_hit = next((mn for mn in tiny_fv_mnemonics if f'"{mn}"' in body), None)
    if has_predicate or mnemonic_hit:
        offenders.append(f"{td_path}:{lineno}: def {name} ({reason})")
\end{lstlisting}
\fcite{llvm/lib/Target/YuShuXin/auto-td/tests/test\_no\_handwritten\_instruction\_td.py:66--112（节选）}

\textbf{守护的实际边界（如实说明）。} 该词法扫描覆盖最直接、最可能的绕过路径——在顶层后端 \texttt{.td} 手写一条 \texttt{def\ ...\ :\ RVInst*} 指令记录。它\textbf{当前不覆盖}：经 \texttt{defm} 多类实例化、经非 \texttt{RVInst} 包装类派生、或置于子目录 \texttt{.td} 的指令记录，也不检查 \texttt{YSXISelDAGToDAG.cpp} 里的手写下降。这些更完整的闭合留给 §9 计划中的 \texttt{llvm-tblgen} 语义门。

\textbf{(2) 边界哨兵}：在 \texttt{YSXInstrInfo.td:1751--1762} 用 \texttt{// YSX-AUTO-TD-BEGIN/END} 围住生成 include，守护断言哨兵区间内不得出现手写 \texttt{def}。\textbf{(3) 数值/往返守护}：断言 \texttt{auto\_full == YAML 数}、\texttt{retained\_schema\_gap == 0}；逐 YAML 校验声明的 pseudo/pattern/builtin 事实都出现在生成 manifest 中。\textbf{(4) 注册门禁}：整套守护既注册为 lit 测试（\texttt{auto-td-guards.test}，失败即让 \texttt{check-llvm}/CI 失败），又注册为构建期目标，使普通 \texttt{ninja} 也门禁：
\begin{lstlisting}
COMMAND ${Python3_EXECUTABLE} -m unittest discover
        -s ${CMAKE_CURRENT_SOURCE_DIR}/auto-td/tests -p "test_*.py"   # :104
add_custom_target(YSXAutoTdGuards ALL DEPENDS ${YSX_AUTO_TD_GUARD_STAMP})  # :118
add_dependencies(YSXCommonTableGen YSXAutoTdGuards)                         # :119
\end{lstlisting}
\fcite{llvm/lib/Target/YuShuXin/CMakeLists.txt:95--119（节选）}

\textbf{验证}：51 个 auto-td 单测全过；向顶层 \texttt{.td} 注入一条手写 tiny-F \texttt{def} 时守护按 \texttt{file:line} 报错并失败，移除后恢复；\texttt{YSXAutoTdGuards} 在构建期执行；覆盖率仍为 \texttt{auto\_full:50}、\texttt{retained\_schema\_gap:0}。

\begin{figure}[ht]
\centering
\begin{tikzpicture}[font=\scriptsize,node distance=4mm and 12mm,
  g/.style={draw,rounded corners=2pt,fill=orange!10,align=center,minimum height=6mm,minimum width=34mm},
  p/.style={draw,fill=blue!6,align=center,minimum height=6mm,minimum width=30mm},
  reg/.style={draw,rounded corners=2pt,fill=green!9,align=center,minimum height=6mm,minimum width=30mm},
  ar/.style={-{Latex[length=1.6mm]}}]
  \node[g] (g1) {源端守护\\\scriptsize test\_no\_handwritten};
  \node[g,below=of g1] (g2) {哨兵 BEGIN/END};
  \node[g,below=of g2] (g3) {数值/往返 + manifest};
  \node[p,right=of g1] (p1) {顶层 .td 指令记录};
  \node[p,right=of g2] (p2) {生成 include 区};
  \node[p,right=of g3] (p3) {coverage + manifests};
  \draw[ar] (g1)--(p1); \draw[ar] (g2)--(p2); \draw[ar] (g3)--(p3);
  \node[reg,below=9mm of g3] (r1) {lit auto-td-guards.test};
  \node[reg,right=of r1] (r2) {CMake YSXAutoTdGuards};
  \node[p,below=4mm of r1,fill=gray!8] (c1) {check-llvm / CI 失败};
  \node[p,below=4mm of r2,fill=gray!8] (c2) {ninja 构建失败};
  \draw[ar] (r1)--(c1); \draw[ar] (r2)--(c2);
  \begin{scope}[on background layer]
    \node[draw,dashed,rounded corners,fit=(g1)(g3)(p1)(p3),inner sep=3mm,label={[font=\scriptsize]above:守护盯产物}] {};
  \end{scope}
\end{tikzpicture}
\caption{守护关系：三类守护各盯一类产物，整套经 lit 与 CMake 双注册，分别门禁 \texttt{check-llvm}/CI 与本地 \texttt{ninja} 构建。}
\label{fig:guard}
\end{figure}

\section{验证与覆盖}
\begin{itemize}[leftmargin=1.4em,itemsep=2pt]
\item 单测：\texttt{unittest discover} $\to$ 51 tests OK。
\item 覆盖（以 \texttt{build/lib/Target/YuShuXin/auto-td/coverage.md} 为权威产物）：\texttt{auto\_full:50}（15 tiny-F + 35 tiny-V）、\texttt{auto\_with\_structured\_override:0}、\texttt{retained\_schema\_gap:0}；其中 8 条带 \texttt{builtin.codegen}（\texttt{validate.py} 白名单恰为这 8 类）具备端到端路径，其余 42 条仅到记录+MC 层。
\item lit（标注所在层）：MC 层 \texttt{llvm/test/MC/YSX/\{tinyf-auto-td.s, tinyv-auto-td.s, tinyv-invalid-disassemble.s\}}；CodeGen 层 \texttt{llvm/test/CodeGen/YSX/\{auto-td-guards.test, tinyv-builtins-isel.ll, tinyf-isel.ll\}}；端到端目标码 \texttt{clang/test/CodeGen/YSX/tinyv-builtins-asm.c}（\texttt{-emit-obj} + \texttt{llvm-objdump} 往返）——均通过。
\item 局限：本轮验证停在生成/文本 + lit 层；生成的 TableGen 已经 lit 工具链消费，但未以 \texttt{llvm-tblgen -print-records} 做语义层门禁（列为 §9 未来工作）。
\end{itemize}

\section{局限与未来工作}
\begin{itemize}[leftmargin=1.4em,itemsep=2pt]
\item \textbf{真实 DAG pattern 生成}：当前 \texttt{*Patterns.inc} 为注释 manifest，全部 8 条 codegen 指令的下降仍是手写 C++；可让生成器产出真正的 VPat 风格模式。
\item \textbf{多 LMUL / 掩码 / 固定类型}：解除 \texttt{v4i32}/\texttt{v4f32} 硬限，按 pseudo matrix 下降。
\item \textbf{tblgen 语义门}：增加一道 \texttt{llvm-tblgen} 守护，断言每个带 tiny-FV 谓词的指令记录名都以 \texttt{YSX\_AUTO\_} 开头——这能闭合 §7 词法守护未覆盖的 \texttt{defm}/包装类/子目录/间接路径。
\item \textbf{InstFormat 细化}：\texttt{\_inst\_format}（\texttt{emit\_td.py:399--407}）当前是粗启发式（向量类恒 \texttt{InstFormatR}），可做精确推断。
\item \textbf{调度面}：taxonomy 的 \texttt{sched:} 当前仅信息性（恒置 \texttt{hasNoSchedulingInfo=1}），可落地。
\item \textbf{从 \texttt{riscv-isa-manual} 自动提炼语义}：现状机器解析仅 \texttt{riscv-opcodes}，\texttt{isa-manual} 当前\textbf{未被任何工具或字段引用}（仅作团队人审旁路参考）。把散文式语义半自动提炼进 taxonomy 是明确未来方向，需 NLP/结构化抽取，\textbf{不在本轮范围}。
\item \textbf{明确不在范围}：完整 RVV 前端类型/可伸缩向量 C ABI、广义自动向量化。
\end{itemize}

\section{结论}
承诺“每条新 tiny-F/tiny-V 指令都从 opcode 源 + 一个 YAML 生成、无手写嵌套 TableGen 指令类”在当前 50 条切片中\textbf{实质达成}：committed \texttt{.td} 无任何 \texttt{YSX\_AUTO\_*} 或 tiny-FV 谓词记录，50 条仅经生成 include 进入构建，每条都是单一共享基类的扁平 \texttt{def}。补齐源端守护与测试注册后，这一承诺在\textbf{守护覆盖的形态内（顶层 \texttt{.td} 的 \texttt{RVInst} 族指令记录）}从约定升级为构建/CI 门禁；\texttt{defm}/包装类/间接路径的完整闭合由计划中的 tblgen 语义门接管。对照表明：声明式数据驱动显著降低单指令的工程与评审成本，把复杂度从“每指令深层间接”迁移到“集中式生成器 + 声明式 taxonomy”；代价是对生成器与数据源正确性的集中依赖、当前仍为手写的下降，以及尚未生成的 pseudo/pattern 矩阵。对“人为知识可由数据派生”的命题，诚实的回答是\textbf{分层成立}：机械层是，设计层否——而把两者分开，正是本工作的价值所在。

\section*{附录 A — 50 条指令的 opcode 来源}
\addcontentsline{toc}{section}{附录 A}
15 tiny-F + 34 tiny-V 来自外部权威 \texttt{riscv-opcodes}（\texttt{rv\_f}/\texttt{rv\_v}）；1 条 \texttt{yushuxin.vfexp} 来自本团队维护的 \texttt{ysx-opcodes}（\texttt{rv\_xtinyv}）。8 条带 \texttt{builtin.codegen}（见 §8）。权威覆盖产物为 \texttt{build/lib/Target/YuShuXin/auto-td/coverage.md}，源真值为 \texttt{auto-td/instructions/\{tiny-f,tiny-v\}/*.yaml} 的文件计数（15 + 35）。

\section*{附录 B — 关键文件/行号索引}
\addcontentsline{toc}{section}{附录 B}
\begin{itemize}[leftmargin=1.4em,itemsep=1pt]
\item 生成器：\texttt{auto-td/tools/ysx\_auto\_td/\{opcodes.py:11--33, loader.py:184--185, validate.py:4--14, emit\_td.py:334--366/399--407, report.py:5--25\}}
\item 数据：\texttt{taxonomy/vector-alu.yaml:66--80}、\texttt{instructions/tiny-v/yushuxin\_vfexp.yaml:1--13}、\texttt{third\_party/ysx-opcodes/extensions/rv\_xtinyv:1}
\item 后端接入与手写部分：\texttt{YSXInstrInfo.td:1751--1762}、\texttt{YSXISelDAGToDAG.cpp:1133--1206}、\texttt{clang/lib/Headers/ysx\_vector.h}
\item 守护：\texttt{test\_no\_handwritten\_instruction\_td.py:66--112}、\texttt{llvm/test/CodeGen/YSX/auto-td-guards.test}、\texttt{CMakeLists.txt:95--119}
\item 端到端测试：\texttt{clang/test/CodeGen/YSX/\{yushuxin-vfexp.c（C$\to$IR）, tinyv-builtins-asm.c（目标码:100）\}}
\item 上游对比：\texttt{RISCVInstrFormatsV.td:107--125}、\texttt{RISCVInstrInfoV.td:612,1137}、\texttt{RISCVInstrInfoVPseudos.td:2160,4810,6111}、\texttt{IntrinsicsRISCV.td:1199,1371}、\texttt{riscv\_vector.td:759}、\texttt{RISCVVEmitter.cpp:407,505,573}
\end{itemize}

\section*{参考文献}
\addcontentsline{toc}{section}{参考文献}
\begin{enumerate}[leftmargin=1.6em,itemsep=1pt]
\item LLVM Project 源码（本仓库 \texttt{llvm/}、\texttt{clang/}）。
\item RISC-V \texttt{riscv-opcodes}（\texttt{third\_party/riscv-opcodes}，外部权威编码数据库）。
\item RISC-V ISA Manual（\texttt{third\_party/riscv-isa-manual}，已 vendored，当前\textbf{未被工具链机器引用}，仅作团队人审旁路参考）。
\item 本仓库设计/计划/实现文档：\texttt{docs/superpowers/\{specs,plans,implementation,blog\}/}。
\end{enumerate}

\end{document}
